%Chargement des paquets
\usepackage{amsmath}
\usepackage{amsfonts}
\usepackage{amssymb}
\usepackage{enumerate}
\usepackage{mathtools}
\usepackage{bbm}
\usepackage{xparse, etoolbox}
\usepackage{enumerate}
\usepackage{enumitem}
\usepackage{mathabx}
\usepackage{babel}[french]

%environment
\usepackage{amsthm}
\usepackage{keytheorems}
\usepackage{hyperref} 

%Interface théorème
\newkeytheorem{lem}[
	name = Lemme
]

\newkeytheorem{prop}[
	name = Proposition
]

\newkeytheorem{thm}[
	name = Théorème
]

\newkeytheorem{coro}[
	name = Corollaire
]

\renewcommand*{\proofname}{Démonstration}

\theoremstyle{definition}
\newtheorem{defn}{Définition}

\theoremstyle{definition}
\newtheorem*{nota}{Notation}

\theoremstyle{definition}
\newtheorem*{limi}{Liminaire}

\theoremstyle{remark}
\newtheorem{rmq}{Remarque}

\theoremstyle{plain}
\newtheorem*{fait}{Fait}


%Ensembles de nombres
\newcommand{\N} {\mathbb{N}}
\newcommand{\Ne}{\N^\ast}
\newcommand{\Z} {\mathbb{Z}}
\newcommand{\Q} {\mathbb{Q}}
\newcommand{\R} {\mathbb{R}}
\newcommand{\Rb}{\overline{\mathbb{R}}}
\newcommand{\Rp}{\R_+}
\newcommand{\Rm}{\R_-}
\newcommand{\K} {\mathbb{K}}
\newcommand{\Ket} {\mathbb{K^{\ast}}}
\newcommand{\Cx}{\mathbb{C}}

%Ensembles, tribus
\newcommand{\A}{\mathbb{A}}
\renewcommand{\P}{\mathcal{P}}
\newcommand{\T}{\mathcal{T}}
\newcommand{\F}{\mathcal{F}}
\newcommand{\C} {\mathcal{C}}
\newcommand{\ouv}{\mathcal{O}}
\newcommand{\B}{\mathcal{B}}
\newcommand{\M}{\mathcal{M}}
\newcommand{\etag}{\mathcal{E}}
\newcommand{\etagOp}{\etag^{+}(\Omega)}
\newcommand{\etagO}{\etag(\Omega)}
%%Lp avant quotientage
\newcommand{\Lic}{\mathcal{L}^1}
\newcommand{\Lpc}{\mathcal{L}^p}
\newcommand{\Linfc}{\mathcal{L}^{\infty}}
\newcommand{\Lifull}{\Lic(\Omega, \B, m)}
\newcommand{\Lpfull}{\Lpc(\Omega, \B, m)}
\newcommand{\Linffull}{\Linfc(\Omega, \B, m)}
%%Lp après quotientage
\newcommand{\Li}{L^1}
\newcommand{\Ld}{L^2}
\newcommand{\Lp}{L^p}
\newcommand{\Linf}{L^{\infty}}
\newcommand{\Listd}{\Li(\Omega, m)} 
\newcommand{\Ldstd}{\Ld(\Omega, m)} 
\newcommand{\Lpstd}{\Lp(\Omega, m)} 

%Quelques opérateurs
\DeclareMathOperator{\codim}{codim}
\DeclareMathOperator{\rg}{rg}
\DeclareMathOperator{\tr}{tr}
\DeclareMathOperator{\vect}{Vect}
\DeclareMathOperator{\ima}{Im}
\DeclareMathOperator{\id}{Id}
\DeclareMathOperator{\sign}{sign}
\DeclareMathOperator{\ext}{ext}
\newcommand{\ind}{\mathbbm{1}}
\newcommand*{\spr}[2]{\left(\,#1\,\middle|\,#2\,\right)}
\newcommand{\interior}[1]{%
  {\kern0pt#1}^{\mathrm{o}}%
}
\DeclareMathOperator{\card}{card}
\DeclareMathOperator{\ppcm}{ppcm}

%Commandes de pure flemme
\newcommand{\veps}{\varepsilon}
\newcommand{\intab}{\int_{a}^{b}}
\newcommand{\intR}{\int_{-\infty}^{\infty}}
\newcommand{\intinf}{\int_{- \infty}^{+ \infty}}
\newcommand{\equi}{\Leftrightarrow}
\newcommand{\Kev}{$\K$-espace vectoriel }
\newcommand{\Kevs}{$\K$-espaces vectoriels }
\newcommand{\Rev}{$\R$-espace vectoriel }
\newcommand{\Revs}{$\R$-espaces vectoriels }
\newcommand{\Cev}{$\Cx$-espace vectoriel }
\newcommand{\Cevs}{$\Cx$-espaces vectoriels }
\newcommand{\evn}{(E, \norm{.})}
\newcommand{\interf}[2]{[#1,#2]}
\newcommand{\limb}{\lim\limits_}
\newcommand{\lebn}{\lambda_n}
\newcommand{\pp}{\text{~presque partout}}
\newcommand{\derivp}[2]{\frac{\partial #1}{\partial #2}}
\newcommand{\famiu}{(u_i)_{i \in I}}
\newcommand{\sev}{\text{~s.e.v.}}
\newcommand{\Mnp}{M_{n,p}(\K)}
\newcommand{\Mn}{M_{n}(\K)}
\newcommand{\tq}{\text{~t.q.~}}
\newcommand{\mq}{~montrer~que~}
\newcommand{\et}{\quad \text{et} \quad}
\newcommand{\ou}{\text{~ou~}}
\newcommand{\Kx}{\K[X]}


%Commandes de gars chelou
\renewcommand{\subset}{\subseteq}

%Commande restriction, ty duckduckgo
\def\restr#1#2{\mathchoice
              {\setbox1\hbox{${\displaystyle #1}_{\scriptstyle #2}$}
              \restraux{#1}{#2}}
              {\setbox1\hbox{${\textstyle #1}_{\scriptstyle #2}$}
              \restraux{#1}{#2}}
              {\setbox1\hbox{${\scriptstyle #1}_{\scriptscriptstyle #2}$}
              \restraux{#1}{#2}}
              {\setbox1\hbox{${\scriptscriptstyle #1}_{\scriptscriptstyle #2}$}
              \restraux{#1}{#2}}}
\def\restraux#1#2{{#1\,\smash{\vrule height .8\ht1 depth .85\dp1}}_{\,#2}} 

%Quelques normes
\DeclarePairedDelimiter\abs{\lvert}{\rvert}
\DeclarePairedDelimiter\labs{\bigl\lvert}{\bigr\rvert}
\DeclarePairedDelimiter\norm{\lVert}{\rVert}
\DeclarePairedDelimiterXPP\normi[1]{}\lVert\rVert{_1}{\ifblank{#1}{\:\cdot\:}{#1}}
\DeclarePairedDelimiterXPP\normd[1]{}\lVert\rVert{_2}{\ifblank{#1}{\:\cdot\:}{#1}}
\DeclarePairedDelimiterXPP\normp[1]{}\lVert\rVert{_p}{\ifblank{#1}{\:\cdot\:}{#1}}
\DeclarePairedDelimiterXPP\normq[1]{}\lVert\rVert{_q}{\ifblank{#1}{\:\cdot\:}{#1}}
\DeclarePairedDelimiterXPP\norminf[1]{}\lVert\rVert{_\infty}{\ifblank{#1}{\:\cdot\:}{#1}}

%Bracket en angle
\DeclarePairedDelimiter\angl{\langle}{\rangle}

%Set, parenthèses
\newcommand{\set}[1]{\{ #1 \}}
\newcommand{\bigset}[1]{\bigl\{ #1 \bigr\}}
\newcommand{\Bigset}[1]{\Bigl\{ #1 \Bigr\}}
\newcommand{\bigpar}[1]{\bigl( #1 \bigr)}
\newcommand{\Bigpar}[1]{\Bigl( #1 \Bigr)}

%Opitmisation des opérateurs de comparaison
\DeclarePairedDelimiter\tio{o (}{)}
\DeclarePairedDelimiter\gro{O (}{)}


\pagestyle{fancy}

\fancyhead[C]{Lemme de Fatou et théorèmes de convergence}
\fancyhead[R]{}

\begin{document}

\underline{Source :} Rudin - Analyse - page 20 \newline

\underline{Leçons :} 
\begin{itemize}
\item 234 : Fonctions et espaces de fonctions Lebesgue-intégrables.
\item 235 : Problèmes d’interversion de symboles en analyse
\item 239 : Fonctions définies par une intégrale dépendant d’un paramètre. Exemples et applications.
\item 241 : Suites et séries de fonctions. Exemples et contre-exemples.
\end{itemize}

\begin{nota}
Dans ce document, $(X, \T, m)$ est un espace mesuré.
Pour toute application $f$ définie sur $X$, on note :
\[ \set{f > a} = \set{x \in X ; f(x) > a} \et \set{f< b} = \set{x \in X ; f(x) < b} \]
\end{nota}

\begin{lem}
Soit $f : X \to \Rb$, les assertions suivantes sont équivalentes :
\begin{enumerate}[label=(\roman*)]
\item $f$ est mesurable ;
\item $\set{f > a} \in \T$ pour tout $a \in \R$ ;
\item $\set{f < b} \in \T$ pour tout $b \in \R$ ;
\end{enumerate}
\end{lem}

\begin{proof}
On rappelle que la tribu borélienne de $\Rb$ est engendré par les intervalles de la forme $]a, \infty]$ ou par les intervalles de la forme
$[-\infty, b[$. De plus, $\set{f > a} = f^{-1}(]a, \infty]$ et $\set{f < b} = f^{-1}([-\infty, b[)$.
\end{proof}

\begin{thm}[de convergence monotone]
Soit $(f_n)_\N$ une suite de fonctions mesurables et positives sur $X$ et supposons que :
\begin{enumerate}[label=(\roman*)]
\item $\forall x \in X, \quad 0 \leq f_0(x) \leq f_1(x) \leq \dots \leq \infty$ ;
\item $\forall x \in X, \quad f_n(x) \to f(x)$ pour $x \to \infty$ ;
\end{enumerate}
autrement dit, la suite $(f_n)_\N$ est croissante et converge simplement vers une fonction $f$.
Alors $f$ est une fonction mesurable et :
\[ \int_X f_n dm \to \int_X f dm \text{ pour } n \to \infty . \]
\end{thm}

\begin{proof}
La suite de réels $\left( \int_X f_n dm \right)_\N$ est croissante, elle est donc convergente dans $[0, \infty]$. 
On note $a$ sa limite. La fonction $f$ est mesurable comme limite simple d'une suite de fonctions mesurables.
De plus, comme $f_n \leq f$, on sait que $\int_X f_n dm \leq \int_X f dm$ pour tout $n \in \N$ et, par passage à la limite :
\[ a \leq \int_X f dm . \]
Soit $s$ une fonction étagée mesurable telle que $0 \leq s \leq f$, et soit $c$ une constante $0<c<1$. On définit :
\[ E_n = \set{x ; f_n(x) \geq cs(x) } \quad \forall n \in \N . \]
Chaque ensemble $E_n = \{ f_n - cs \geq 0 \}$ est mesurable, $E_0 \subset E_1 \subset \dots$ et $X = \bigcup E_n$.
En effet, si $x=0$ alors $x \in E_0$, sinon $f(x) >0, cs(x) < f(x)$ (car $c<1$) et par définition de la limite, il existe un rang à partir duquel
$f_n(x) \geq cs(x)$, donc $x \in E_n$ pour au moins un rang $n$. Par conséquent :
\[ \int_X f_n dm \geq \int_{E_n} f_n dm \geq c \int_{E_n} s dm \quad \forall n \in \N . \]
L'application $E \mapsto \int_{E} s dm$ définit une mesure de $\T$ et, par $\sigma$-additivité :
\[ \lim_{n \to \infty} \int_{E_n} s dm = \int_X s dm , \]
et donc, par passage à la limite :
\[ a \geq c \int_X s dm . \]
Ceci étant vrai pour toute constante $c \in ]0,1[$ et pour toute fonction étagée mesurable telle que $0 \leq s \leq f$, on déduit :
\[ a \geq \inf_X f dm . \]
Finalement, $a = \inf_X f dm$, ce qui démontre le résultat.
\end{proof}

\begin{lem}[de Fatou]
Soit $(f_n)_\N$ une suite de fonctions mesurables positives définies sur $X$. On a :
\[ \int_X \liminf f_n dm \leq \liminf \int_X f_n dm . \] 
\end{lem}

\begin{proof}
On définit :
\[ g_k(x) = \inf_{i \geq k} f_i(x) \quad \forall k \geq 1 ; \forall x \in X , \]
de sorte que :
\[ \int_X g_k dm \leq \inf_{i \geq k} \int_X f_i dm \quad \forall k \geq 1 . \]
Par ailleurs, la suite $(g_k)_\N$ est croissante et, la mesurabilité étant stable par passage à la borne inférieure, on peut appliquer 
le théorème de convergence monotone à la suite $(g_k)_\N$, ce qui donne :
\begin{align*}
\int_X \liminf f_n dm & = \int_X \lim_{n \to \infty} g_n dm \\
& = \lim_{n \to \infty} \int_X g_n dm \\
& \leq \lim_{n \to \infty} \inf_{i \geq n} \int_X f_i dm \\
& \leq \liminf \int_X f_n dm
\end{align*}
\end{proof}

\begin{thm}[de convergence dominée]
Soit $(f_n)_\N$ une suite de fonctions complexes mesurables sur $X$. On définit $f$ sur $X$ par :
\[ f(x) = \lim_{n \to \infty} f_n(x) . \]
On suppose qu'il existe une fonction $g$ d'intégrale finie et telle que :
\[ \abs{f_n(x)} \leq g(x) \quad \forall n \geq 1 ; \forall x \in X , \]
alors l'intégrale de $f$ est finie, et :
\begin{align} \lim_{n \to \infty} \int_X \abs{f_n-f} dm = 0 , \end{align}
ainsi que :
\begin{align} \lim_{n \to \infty} \int_X f_n dm = \int_X f dm . \end{align}
\end{thm}

\begin{proof}
Par passage à la limite, on a $\abs{f} \leq g$ sur $X$, donc l'intégrale de $f$ est finie.
On remarque que $2g - \abs{f_n - f} \geq 0$, on applique le lemme de Fatou à cette suite :
\begin{align*}
\int_X 2g dm & = \int_X \liminf (2g - \abs{f_n - f}) dm \\
& \leq \liminf \int_X (2g - \abs{f_n - f}) dm \\
& \leq \int_X 2g dm + \liminf \left ( - \int_X \abs{f_n - f}) dm \right ) \\
& \leq \int_X 2g dm - \limsup \int_X \abs{f_n - f} dm .
\end{align*}
Puisque $\int_X 2g dm$ a une valeur finie, on le retranche des deux membres de l'inégalité pour obtenir :
\[ \limsup \int_X \abs{f_n - f}) dm \leq 0 , \]
Ce qui implique nécessairement (1). 
Par ailleurs, comme $\abs*{\int_X f dm} \leq \int_X \abs{f} dm$, l'inégalité (1) implique :
\[ 0 \leq \lim_{n \to \infty} \abs*{ \int_X f_n dm - \int_X f dm} \leq 0 , \]
ce qui implique (2).
\end{proof}

\end{document}